\color{uniblue}\section[Описание и работа]{Описание и работа\color{white!0}<ABBRS>} % вся эта белиберда только чтоб невидимый тэг воткнуть!
\color{black}

\color{uniblue}{\subsection{Назначение}}
\color{black}

\begin{enumerate}[label=\arabic{section}.\arabic{subsection}.\arabic*, labelsep=4pt, leftmargin=0pt, itemindent=57pt]

\item
Устройство серии <<ЮНИТ-М300-АРНТ>> предназначено  для  оперативного  и  автоматического  управления приводом РПН силового трансформатора с номинальным значением стороны высшего напряжения 6-220 кВ с целью поддержания определенного уровня напряжения на шинах подстанции или у потребителя при  изменениях  режима  работы  сети.  Устройство  поддерживает  режим  регулирования  двух трансформаторов, включенных на параллельную работу. Устройство может применяться для управления приводом устройств РПН двухобмоточных (включая ЛРТ),  трехобмоточных  трансформаторов,  а  также  двухобмоточных  трансформаторов  с  расщепленной обмоткой.
\item
Устройство может осуществлять функции измерения, регистрации, осциллографирования и сигнализации. При подключении устройства к информационной сети энергообъекта возможно осуществление функций дистанционного управления устройством и сбор данных с него посредством протоколов и интерфейсов, определяемыми конфигурацией устройства.
\item 
Устройство предлагается с типовой конфигурацией: в технической части определяемой картой заказа (см. приложение \ref{app:card}), в функциональной части -- функционально-структурной схемой (см. приложение \ref{app:fsu}). Схемы подключения устройства по аналоговым входам приведены в приложении \ref{app:analogs}, схемы типового подключения сигналов алгоритмов к дискретным входам и релейным выходам устройства приведены в приложениях \ref{app:inouts}, \ref{app:inouts2}.
\item
По желанию Заказчика возможно выполнение устройства в других конструктивных исполнениях в соответствии с общей картой заказа на серию <<ЮНИТ-М300>>, которая приводится в документации ЮТКБ.656122.609 РЭ <<Устройство защиты и автоматики серии <<ЮНИТ-М300>>. Руководство по эксплуатации общее>>. 
\item
Конструктивное исполнение и функциональный состав Устройство в базовом исполнении предлагается в двух версиях: \begin{itemize} 
\item 
для ПС I и смешанной архитектуры (далее по тексту – исполнение I); 
\item 
для ПС II архитектуры (далее по тексту – исполнение II), \end{itemize} в аппаратной части определяемой картой заказа устройства в корпусе М314 (см. приложение Б), в функциональной части – структурно-функциональной схемой (см. приложение Г). Подключение устройства по аналоговым входам приведено в приложении Д , типовое подключение сигналов алгоритмов к дискретным входам и релейным выходам устройства приведено в приложении Е.
\item
Устройство серии <<ЮНИТ-М300-АРНТ>> может применяться в составе шкафов для установки на объектах филиалов и ДО ПАО <<Россети>>, построенных на базе архитектур I и II типа, а также на базе смешанной архитектуры. 

\end{enumerate}

%==================== Технические характеристики устройства =============================

\color{uniblue}\subsection{Технические характеристики}
\color{black}

\begin{enumerate}[label=\arabic{section}.\arabic{subsection}.\arabic*, labelsep=4pt, leftmargin=0pt, itemindent=57pt]

\item
Основные номинальные параметры устройства указаны в таблице 1.1.\ref{desc:chars}. 

\setlength{\extrarowheight}{0.06cm} % высота ячеек в таблице
\small
\begin{longtable}
    {|>{\centering\arraybackslash}m{12.9cm}|>{\centering\arraybackslash}m{4cm}|}
    \caption{Основные номинальные параметры устройства\hfill\vspace{-0.5\baselineskip}}\label{desc:chars}\\ 
    \hhline{|-|-|}
    \rowcolor{gray!30}
    Параметр & Значение \\
    \hhline{|-|-|}
    \endfirsthead

    %\caption{\emph{Продолжение\hfill\vspace{-0.5\baselineskip}}} \\ % Отступ обозначения таблицы от таблицы и выравнивание надписи слева
    \caption*{\hspace{3pt}\emph{Продолжение таблицы \ref{desc:chars}\hfill\vspace{-0.5\baselineskip}}} \\ % сделано по ГОСТ 2.105 п.6.8.7
    \hhline{|-|-|}
    \rowcolor{gray!30}
    Параметр & Значение \\ 
    \hhline{|-|-|}
    \endhead
    \hhline{|-|-|}

    \endfoot
    \endlastfoot

    \raggedright Номинальный ток, А  & \raggedright\arraybackslash 1; 5 \\  \hhline{|-|-|}
    \raggedright Номинальный переменное фазное напряжение, В  & \raggedright\arraybackslash 100/$\sqrt{3}$ \\  \hhline{|-|-|}
    \raggedright Номинальная частота, Гц  & \raggedright\arraybackslash 50 \\  \hhline{|-|-|}
    \raggedright Номинальное напряжение оперативного постоянного тока, В  & \raggedright\arraybackslash 110; 220 \\  \hhline{|-|-|}
    \raggedright Номинальное напряжение оперативного переменного тока, В  & \raggedright\arraybackslash 110; 220 \\  \hhline{|-|-|}

\end{longtable}
\normalsize
\setlength{\extrarowheight}{0.0cm} % возврат параметра высота ячеек в таблице

\end{enumerate}
%======================================= Состав и конструктивное исполнение ===============================================

\color{uniblue}\subsection{Состав и конструктивное исполнение}
\color{black}

\begin{enumerate}[label=\arabic{section}.\arabic{subsection}.\arabic*, labelsep=4pt, leftmargin=0pt, itemindent=57pt]

\item
Корпус устройства в типовой конфигурации предусматривает отдельные исполнения по ширине в зависимости от требуемого количества модулей дискретного управления и дискретных входов, габаритные размеры устройства приведены в приложении \ref{app:sizes}. Внешний вид устройства для различных типовых исполнений показан на рисунках \ref{device:img1}, \ref{device:img2}. Подробное описание конструкции и функций используемых модулей (а также доступных к заказу) приведено в документации ЮТКБ.656122.609 РЭ <<Устройство защиты и автоматики серии <<ЮНИТ-М300>>. Руководство по эксплуатации общее>>.

\vspace{3mm}
\begin{figure}[H]
\centering
\includegraphics[width=0.7\textwidth,height=0.7\textheight,keepaspectratio]{314.jpg}
\captionof{figure}{ Вид устройства ЮНИТ-М314-АРНТ для ПС со смешанной архитектурой с размещением плат 
по слотам (исполнение I)}\label{device:img1}
\end{figure}

\vspace{3mm}
\begin{figure}[H]
\centering
\includegraphics[width=0.9\textwidth,height=0.9\textheight,keepaspectratio]{319.jpg}
\captionof{figure}{- Вид устройства ЮНИТ-М314-АРНТ для ПС со второй архитектуры с размещением плат по 
слотам (исполнение II)}\label{device:img2}
\end{figure}

\item
Перечень функций, реализованных в составе устройства, приведен в таблице \ref{desc:tbl1}. Описание основных функций устройства приводится в разделе \ref{sec:funcs} данной документации. Подробное описание вспомогательных (сервисных) функций приведено в документации ЮТКБ.656122.609 РЭ <<Устройство защиты и автоматики серии <<ЮНИТ-М300>>. Руководство по эксплуатации общее>>.

\setlength{\extrarowheight}{0.05cm} % высота ячеек в таблице
\small
\begin{longtable}
    {|>{\centering\arraybackslash}m{3.2cm}|>{\centering\arraybackslash}m{13.7cm}|}
    \caption{Функции в составе устройства <<ЮНИТ-М300-АРНТ>>\hfill\vspace{-0.5\baselineskip}}\label{desc:tbl1}\\ 
    \hhline{|-|-|}
    \rowcolor{gray!30}
    Группа функций & Примечание \\ 
    \hhline{|-|-|}
    \endfirsthead

    \caption*{\hspace{3pt}\emph{Продолжение таблицы \ref{desc:tbl1}\hfill\vspace{-0.5\baselineskip}}} \\ % сделано по ГОСТ 2.105 п.6.8.7
    \hhline{|-|-|}
    \rowcolor{gray!30}
    Группа функций & Примечание \\ 
    \hhline{|-|-|}
    \endhead
    \hhline{|-|-|}

    \endfoot
    \endlastfoot

    \centering Основные функции & \flushleft\arraybackslash Управление приводом РПН: \begin{itemize} 
    \item   ручное регулирование напряжения;
    \item   автоматическое поддержание уровня напряжения на шинах или у потребителя;
    \item  встречное регулирование (коррекция уровня напряжения на  шинах у потребителя в зависимости от тока нагрузки);
    \item  оперативный выбор уровня поддерживаемого напряжения;
    \item  одновременный контроль двух секций с возможностью автоматического переключения между секциями;
    \item   групповое регулирование двух параллельно работающих  трансформаторов;
    \item  управление электроприводом импульсными или непрерывными  командами;
    \item   контроль приводного механизма в процессе переключения;
    \item  блокировки по температуре и уровню масла;
    \item  блокировки при перегрузках;
    \item  блокировки по повышению/понижению, несимметрии напряжения;
    \item   блокировки при обнаружении неисправности привода;
    \item  контроль механического ресурса привода.
\end{itemize} \\  \hline

    \multirow{7}{*}{\centering Сервисные функции} & \flushleft\arraybackslash Сигнализация: \begin{itemize}
    \item обеспечение ввода служебных дискретных сигналов о состоянии оборудования шкафа РЗА;
    \item формирование отчетов по изменению состояния дискретных внешних и внутренних сигналов;
    \item  выдача отчетов в АСУ ТП;
    \item формирование и выдачу сигналов аварийной и предупредительной сигнализации в схему резервной сигнализации;
    \end{itemize} \\ \cline{2-2}
    & \flushleft\arraybackslash  Осциллографирование и журналирование:
    \begin{itemize} 
    \item 
    запись доаварийных, аварийных и послеаварийных режимов с сохранением осциллограмм в энергонезависимой памяти устройства;
    \item 
    запись о событиях по изменению конфигурации устройства, изменения состояния внутренних сигналов алгоритмов, обновление встраиваемого программного обеспечения и пр. с метками времени и описанием события в энергонезависимой памяти устройства.
    \end{itemize}
    \\ \cline{2-2}
    & \flushleft\arraybackslash Самодиагностика: 
    \begin{itemize}
    \item 
    непрерывный оперативный контроль работоспособности в течение всего времени работ.
    \end{itemize} \\ \cline{2-2}
    & \flushleft\arraybackslash Конфигурирование: 
    \begin{itemize}
    \item задание/переключение уставок, ввод в работу / вывод из работы;
    \item изменение параметров самоописания устройства;
    \item перезагрузка устройства;
    \item ввод/вывод тестового режима; 
    \item изменение режима работы информационных сервисов; 
    \item возможность переопределения связей входных и выходных логических сигналов алгоритмов устройства с дискретными входами, выходными реле, светодиодными индикаторами и функциональными клавишами управления.
    \end{itemize} \\ \cline{2-2}
    & \flushleft\arraybackslash Тестовый режим: 
     \begin{itemize}
     \item
    обработка принятых сигналов и отправка отчетов с установленным флагом <<Тест>>.
    \end{itemize} \\ \cline{2-2}
    & \flushleft\arraybackslash Самоописание: 
    \begin{itemize}
    \item
    предоставление данных об устройстве посредством устройства <<ЮНИТ-ИЧМ>> и сервисную программное обеспечение <<ЮНИТ-СЕРВИС>>
    \end{itemize} \\ \cline{2-2}
    & \flushleft\arraybackslash Синхронизация времени: 
    \begin{itemize} 
    \item
    календарь и часы реального времени, синхронизируемые от персонального компьютера или от АСУ~ТП 
    \end{itemize} \\ \hhline{|-|-|}

\end{longtable}
\setlength{\extrarowheight}{0.0cm} % возврат параметра высота ячеек в таблице
\normalsize

\end{enumerate}

%============================================================ Устройство и работа ==================================================
\color{uniblue}\subsection{Устройство и работа}\color{black}

\begin{enumerate}[label=\arabic{section}.\arabic{subsection}.\arabic*, labelsep=4pt, leftmargin=0pt, itemindent=57pt]
\item
Описание функций релейной защиты и автоматики в составе устройства, а также их функциональные схемы и уставки приводятся в разделе \ref{sec:funcs}. Общая функционально-структурная схема взаимосвязей функций, дискретных входов и выходов устройства приведена в приложении \ref{app:fsu}. Принятые обозначения элементов схем приведены в приложении \ref{app:func_desc}. %в исходнике приложения Г и Ж соответственно
\item
Система самодиагностики устройства выполняет тесты в полном объеме при включении устройства, а также непрерывно в фоновом режиме.
Системой самодиагностики устройства в автоматическом режиме контролируются следующие параметры:
\begin{itemize}
\item состояние аппаратной части устройства, в том числе: АЦП модулей ввода аналоговых сигналов, блока питания, ОЗУ, ПЗУ, процессора, модулей дискретных входов, модулей дискретного управления, вспомогательных модулей;
\item температурный режим устройства;
\item состояние синхронизации времени;
\item сохранность исполняемого программного кода (целостность ПО);
\item состояние измерительных цепей;
\item исправность используемых каналов связи.
\end{itemize}
\item
Система самодиагностики фиксирует критические неисправности (<<Неисправность>>) и некритические ошибки (<<Ошибка>>). В случае появления критической неисправности активируется внутренний сигнал <<Критическая неисправность>>, замыкаются контакты реле IRF, прекращается работа алгоритмов основных функций в составе устройства, это состояние устройства сигнализируется свечением светодиода <<Неисправность>> и отсутствием свечения светодиода <<Готовность>>. В случае появления некритической ошибки активируется внутренний сигнал <<Некритическая ошибка>>, алгоритмы основных функций в составе устройства продолжают работать, это состояние устройства сигнализируется свечением светодиода <<Неисправность>> и отсутствием свечения светодиода <<Готовность>>. Подробное описание системы самодиагностики приводится в документации ЮТКБ.656122.609 РЭ <<Устройство защиты и автоматики серии <<ЮНИТ-М300>>. Руководство по эксплуатации общее>>.
\item
В данном устройстве предусматривается установка платы аналоговых входов типа <<М046>>
устройства М314 в слоты <<М08>>, <<М09>>. В исполнении II дополнительно предусматривается установка 
платы аналоговых входов типа <<М030>> в слоты <<М06>>, <<М07>>. Типовое подключение входных аналоговых 
сигналов к плате аналоговых входов в составе устройства показано на рисунке Д.1 приложения Д .
Подробное описание параметров настроек и технических характеристик платы аналоговых входов данного 
типа приводится в документации ЮТКБ.656122.609 РЭ <<Устройство защиты и автоматики серии <<ЮНИТ-М300>>. Руководство по эксплуатации общее»
\item
В данном устройстве предусматривается установка плат дискретных входов типа <<В021>> в слоты
<<М5, М6>> в исполнения I (в исполнении II только в слот <<М3>>).Типовое подключение входных сигналов 
функций к дискретным входам в составе устройства показано на рисунках Е.1, Е.2 в приложении Е.
Подробное описание параметров настроек и технических характеристик плат дискретных входов данных 
типов приводится в документации ЮТКБ.656122.609 РЭ <<Устройство защиты и автоматики серии <<ЮНИТ-М300>>. Руководство по эксплуатации общее>>.
\item
В устройствах с поддержкой протокола передачи данных МЭК 61850 в качестве входных сигналов ФСУ могут быть использованы данные из GOOSE-сообщений, на которые оформлена подписка. Количество сигналов GOOSE, получаемых от интеллектуальных устройств, в данном типоисполнении устройства \textcolor{red}{равно 8 шт}.
\item
Сигналы из GOOSE-сообщений проходят предварительную обработку перед тем как попасть в <<Блок объединительный>> и далее в ФСУ (см. приложение \ref{app:fsu}).
В предварительную обработку входят следующие операции:
\begin{itemize}
\item проверка качества принимаемого сигнала;
\item алгоритм обработки резервирования GOOSE-сообщений;
\item алгоритм обработки симуляции GOOSE-сообщений;
\item алгоритм обработки метки <<Test>> GOOSE-сообщений;
\item подстановка значения по умолчанию при <<плохом>> качестве данных, принимаемых из GOOSE-сообщений.
\end{itemize}
Подробнее о подписке на GOOSE-сообщения указано в RU.37182817.00001.02.34.01 – <<Программное обеспечение мониторинга и параметрирования <<ЮНИТ-СЕРВИС>>. Руководство пользователя>>.

\item 
В устройствах с поддержкой протокола передачи данных МЭК 61850 для связи со смежными устройствами может быть использован протокол GOOSE по МЭК 61850-8.1. Количество сигналов GOOSE-сообщений, передаваемых в смежные интеллектуальные устройства, зависит от типоисполнения ЮНИТ-М300, первичного оборудования и в общем случае составляет не менее 100 шт. Подробнее о публикации GOOSE-сообщений указано в RU.37182817.00001.02.34.01 – <<Программное обеспечение мониторинга и параметрирования <<ЮНИТ-СЕРВИС>>. Руководство пользователя>>.
\item
Устройство поддерживает два режима оперативного управления: местный и дистанционный. Выбор режима управления устройством производится посредством функциональной клавиши <<М/Д>> или по состоянию выбранного дискретного входа. В местном режиме игнорируются все команды управления с верхнего уровня АСУ ТП. Более подробно режимы устройства описаны в документации ЮТКБ.656122.609 РЭ <<Устройство защиты и автоматики серии <<ЮНИТ-М300>>. Руководство по эксплуатации общее>>.
\item
Для управления режимами работы основных функций (оперативный ввод/вывод, перевод на сигнал и пр.), а также для генерации различных сигналов сброса функций (сброс триггеров фиксации неисправности цепей газовых и технологических защит, сброс счетчиков и пр.), применяются виртуальные ключи и кнопки. Подробное описание и сценарии работы виртуальных ключей и кнопок приведено в документации ЮТКБ.656122.609 РЭ <<Устройство защиты и автоматики серии <<ЮНИТ-М300>>. Руководство по эксплуатации общее>>. 
\item
В данном типоисполнении устройства реализовано четыре группы уставок. Подробное описание работы с группами уставок приведено в документации ЮТКБ.656122.609 РЭ <<Устройство защиты и автоматики серии <<ЮНИТ-М300>>. Руководство по эксплуатации общее>>.
\item
Устройство предусматривает возможность синхронизации внутренних часов реального времени несколькими способами:
\begin{itemize}
    \item по протоколу SNTP;
    \item по протоколу PTP;
    \item через дискретный вход устройства с использованием сигнала синхронизации 1PPS;
    \item {средствами стандартизованного протокола передачи данных.}
\end{itemize}
Подробное описание способов синхронизации времени устройства приведено в документации ЮТКБ.656122.609 РЭ <<Устройство защиты и автоматики серии <<ЮНИТ-М300>>. Руководство по эксплуатации общее>>.
\item
В данном устройстве предусматривается установка плат выходных реле типа <<К002>> и <<К001>> 
соответственно в слоты <<М3>> и <<М4>> в исполнении I (в исполнении II установка плат выходных реле не 
предусмотрена). Типовое подключение выходных сигналов функций к выходным реле в составе 
устройства показано на рисунке Е.1 приложения Е. Подробное описание параметров настроек и 
технических характеристик плат выходных реле данного типа приводится в документации 
ЮТКБ.656122.609 РЭ <<Устройство защиты и автоматики серии <<ЮНИТ-М300>>. Руководство по 
эксплуатации общее>>.
\item
В устройствах с поддержкой протокола передачи данных МЭК 61850 для связи со смежными 
устройствами может быть использован протокол GOOSE по МЭК 61850-8.1. Количество сигналов GOOSE 
сообщений, передаваемых в смежные интеллектуальные устройства, зависит от типоисполнения <<ЮНИТ-М300>>, первичного оборудования и в общем случае составляет не менее 100 шт. Подробнее о публикации 
GOOSE сообщений указано в RU.37182817.00001.02.34.01 – <<Программное обеспечение мониторинга и 
параметрирования <<ЮНИТ-СЕРВИС>>. Руководство пользователя>>.
\item
Устройство поддерживает два режима оперативного управления: местный и дистанционный. Выбор
режима управления устройством производится посредством функциональной клавиши <<М/Д>> или по 
состоянию выбранного дискретного входа. В местном режиме игнорируются все команды управления с верхнего уровня АСУ ТП. Более подробно режимы устройства описаны в документации ЮТКБ.656122.609 
РЭ <<Устройство защиты и автоматики серии <<ЮНИТ-М300>>. Руководство по эксплуатации общее>>.
\item
Для управления режимами работы основных функций (оперативный ввод/вывод и пр.), а также для 
генерации различных сигналов сброса функций (сброс счетчиков переключений и пр.), применяются так 
называемые виртуальные ключи и кнопки. Подробно описание и сценарии работы виртуальных ключей и 
кнопок приведено в документации ЮТКБ.656122.609 РЭ <<Устройство защиты и автоматики серии <<ЮНИТ-М300>>. Руководство по эксплуатации общее>>.
\item
В устройстве реализовано четыре группы уставок. Подробное описание работы с группами уставок 
приведено в документации ЮТКБ.656122.609 РЭ <<Устройство защиты и автоматики серии <<ЮНИТ-М300>>. 
Руководство по эксплуатации общее>>.
\par
\item
Устройство предусматривает возможность синхронизации внутренних часов реального времени
несколькими способами:
\begin{itemize}
\item по протоколу SNTP;
\item по протоколу PTP;
\item через дискретный вход устройства с использованием сигнала синхронизации 1PPS;
\item средствами стандартизованного протокола передачи данных.
\end{itemize}
\item
Подробное описание способов синхронизации времени устройства приведено в документации 
ЮТКБ.656122.609 РЭ <<Устройство защиты и автоматики серии <<ЮНИТ-М300>>. Руководство по 
эксплуатации общее>>.
\item
Устройство предусматривает возможность ведения журнала событий. Перечень сигналов, 
доступных для регистрации в журнале событий, приведен в приложении А. События безопасности 
записываются в отдельный журнал безопасности. Подробное описание параметров журналов событий 
приведено в документации ЮТКБ.656122.609 РЭ <<Устройство защиты и автоматики серии <<ЮНИТ-М300>>. 
Руководство по эксплуатации общее>>.
\item
Устройство предусматривает возможность осциллографирования событий. Перечень дискретных 
сигналов, доступных для регистрации аварийным осциллографом, приведен в приложении А. Подробное 
описание параметров встроенного регистратора аварийных событий приведено в документации 
ЮТКБ.656122.609 РЭ <<Устройство защиты и автоматики серии <<ЮНИТ-М300>>. Руководство по 
эксплуатации общее>>.
\setlength{\extrarowheight}{0.05cm}
\small
\begin{longtable}{|>{\centering\arraybackslash}m{6.8cm}|>{\centering\arraybackslash}m{2cm}|>{\centering\arraybackslash}m{5.3cm}|>{\centering\arraybackslash}m{2cm}|}
\caption{Перечень аналоговых сигналов, доступных для регистрации аварийным осциллографом\hfill\vspace{-0.5\baselineskip}}\label{osz:tbl1}\\ 
\hline
\rowcolor{gray!30}
Описание аналогового сигнала & Канал & Наименование  & Пуск РАС \\ 
\hline
\endfirsthead
\caption*{\hspace{3pt}\emph{Продолжение таблицы \ref{osz:tbl1}\hfill\vspace{-0.5\baselineskip}}} \\ % сделано по ГОСТ 2.105 п.6.8.7
\hline
\rowcolor{gray!30}
Описание аналогового сигнала & Канал & Наименование  & Пуск РАС \\ 
%\hline
\endhead
%\hline
\endfoot
%\hline
\endlastfoot

\raggedright Напряжение фазы A (междуфазное AB) секции 1 & \centering 1 & \centering U1  & \centering\arraybackslash - \\
\hline
\raggedright Напряжение фазы B (междуфазное BC) секции 2 & \centering 2 & \centering U2  & \centering\arraybackslash - \\
\hline
\raggedright Напряжение фазы C («разомкнутого треугольника») секции 1 & \centering 3 & \centering U3  & \centering\arraybackslash - \\
\hline

\raggedright Напряжение фазы A (междуфазное AB) секции 2 & \centering 4 & \centering U4  & \centering\arraybackslash - \\
\hline
\raggedright Напряжение фазы B (междуфазное BC) секции 2              & \centering 5 & \centering U5  & \centering\arraybackslash - \\
\hline
\raggedright Напряжение фазы C («разомкнутого треугольника») секции 2               & \centering 6 & \centering U6  & \centering\arraybackslash - \\
\hline

\raggedright Ток фазы вводного выключателя 1 секции  & \centering 7 & \centering I1  & \centering\arraybackslash - \\
\hline

\raggedright Ток фазы секционного выключателя 1 секции  & \centering 8 & \centering I2  & \centering\arraybackslash - \\
\hline

\raggedright Ток фазы вводного выключателя 2 секции & \centering 9 & \centering I3  & \centering\arraybackslash - \\
\hline
\raggedright Ток фазы секционного выключателя 2 секции  & \centering 10 & \centering I4  & \centering\arraybackslash - \\
\hline

\raggedright Ток фазы стороны трансформатора с РПН
(только для исполнения II) & \centering 11 & \centering I5  & \centering\arraybackslash - \\
\hline
\raggedright Токовая петля 1 логометра РПН
(только для исполнения I) & \centering 12 & \centering I6  & \centering\arraybackslash - \\
\hline
\raggedright Токовая петля 2 логометра РПН параллельно работающего тр-ра
(только для исполнения I) & \centering 13 & \centering I7  & \centering\arraybackslash - \\
\hline
\raggedright Вычисленное напряжение фазы A
секции 1 & \centering 14 & \centering U7  & \centering\arraybackslash - \\
\hline
\raggedright Вычисленное напряжение фазы B
секции 1  & \centering 15 & \centering U8  & \centering\arraybackslash - \\
\hline
\raggedright Вычисленное напряжение фазы C
секции 1 & \centering 16 & \centering U9  & \centering\arraybackslash - \\
\hline
\raggedright Вычисленное междуфазное напряжение AB секции 1 & \centering 17 & \centering U10  & \centering\arraybackslash - \\
\hline
\raggedright Вычисленное междуфазное напряжение BC секции 1 & \centering 18 & \centering U11  & \centering\arraybackslash - \\
\hline
\raggedright Вычисленное междуфазное напряжение CA секции 1 & \centering 19 & \centering U12  & \centering\arraybackslash - \\
\hline
\end{longtable}
\normalsize
\setlength{\extrarowheight}{0.0cm}

\item В устройстве реализована поддержка режима тестирования <<Тест>> в объеме требований стандартов МЭК 61850. Особенности перевода устройства в режим тестирования и функциональности устройства в этом режиме описаны в документации ЮТКБ.656122.609 РЭ <<Устройство защиты и автоматики серии <<ЮНИТ-М300>>. Руководство по эксплуатации общее>>.

\item\sloppy Описание и настройки доступных портов и интерфейсов связи приводится в документации ЮТКБ.656122.609 РЭ <<Устройство защиты и автоматики серии <<ЮНИТ-М300>>. Руководство по эксплуатации общее>>.

\item Устройство имеет встроенную реализацию правил разграничения прав доступа, для этого в устройстве существует ролевая модель разграничения прав доступа.
В устройстве предусмотрены следующие роли:
\begin{itemize}
    \item Гость (пароль не требуется, доступен просмотр данных устройства);
    \item Дежурный;
    \item Инженер;
    \item Офицер (сотрудник) безопасности;
    \item Администратор (пароль <<АААА>> (кириллица), учетная запись предназначена для первоначального конфигурирования штатного профиля пользователей).
\end{itemize}
Все пользователи делятся на непривилегированных (роль Гость) и привилегированных (все остальные роли).
\item
Предустановленные пользователи с именами Гость и Администратор не подлежат удалению и принудительной блокировке, их роли не подлежат изменению. Гость используется как роль по умолчанию для доступа к просмотру данных устройства в режиме чтение без запроса аутентификационных данных пользователя. Администратор необходим для добавления/удаления других пользователей, установки им соответствующих ролей и установке/изменении паролей пользователей.
\item
После перезагрузки устройства подключенная сессия привилегированного пользователя сбрасывается. Устройство запускается с сеансом непривилегированного пользователя (роль Гость).
Подробное описание ролей и доступных полномочий и ограничений для каждой из ролей приведено в документации ЮТКБ.656122.609 РЭ <<Устройство защиты и автоматики серии <<ЮНИТ-М300>>. Руководство по эксплуатации общее>>.

\end{enumerate}

%============================================================ Средства измерения, инструмент и принадлежности ==================================================
\color{uniblue}\subsection{Средства измерения, инструмент и принадлежности}\color{black}
Перечень оборудования и средств измерения, рекомендуемый для проведения эксплуатационных проверок устройства, приведен в документации ЮТКБ.656122.609 РЭ <<Устройство защиты и автоматики серии <<ЮНИТ-М300>>. Руководство по эксплуатации общее>>.

%================================================================= Маркировка и пломбирование ==================================================
\color{uniblue}\subsection{Маркировка и пломбирование}\color{black}
Сведения по маркировке и пломбировании устройства приведены в документации ЮТКБ.656122.609 РЭ <<Устройство защиты и автоматики серии <<ЮНИТ-М300>>. Руководство по эксплуатации общее>>.

%======================================================================== Упаковка ==================================================
\color{uniblue}\subsection{Упаковка}\color{black}
Сведения по упаковке устройства приведены в документации ЮТКБ.656122.609 РЭ <<Устройство защиты и автоматики серии <<ЮНИТ-М300>>. Руководство по эксплуатации общее>>.
